% =====================================================================
%  Harald Høffding, Epilog
%  In: Tilskueren. Maanedsskrift for Literatur, Samfundsspørgsmaal og almenfattelige videnskabelige Skildringer,
%  udgivet af N. Neergaard.  Syvende Aargang.  Kjøbenhavn: P. G. Philipsens Forlag; Thieles Bogtrykkeri.  Maj 1890, pp. 343-353.
%  Antiqua.
%
%  TRANSCRIPTION -- GENERATED; DO NOT EDIT BY HAND.  Høffding's four pieces in Tilskueren 1886-1890 (Om Religionsfrihed;
%  the three contributions to the Nietzsche debate with Georg Brandes) are built by ONE pipeline
%  (ebook-method/tilskueren-1886-90/) into a combined working file, which split.py cuts into four transcriptions.
%  Rebuild from the repo root:   sh ebook-method/tilskueren-1886-90/rebuild.sh
%
%  STATUS: COMPLETE (2026-10-15): the volume's title page and the whole text, pp. 343-353, with its notes.
%
%  EDITION.  The only printing in Tilskueren.  Text of record: Google's scan of the Harvard copy, Internet Archive
%  tilskuerenmaaned1890nngoog, 600 ppi grey JP2 leaves with the Archive's ABBYY OCR; extract ~/bibliotek/Høffding, Harald/1890-tilskueren-7-pp125-153-343-353.pdf (SCANS.tsv, sha256),
%  made by ia_extract.py (title leaf, Indhold leaves, the article).  Page range confirmed from the volume's Indhold and the
%  signature.  Høffding's last word after Brandes's »Duplik« (Tilskueren, April 1890).  Translated by William Banks, The Great Debate (2024); not translated here.
%
%  ACCURACY (2026-10-15; one collation for all four texts, 76 pp., ~29,000 words).
%  (a) The DRAFT is the Archive's ABBYY OCR (the layer of the working PDF), turned into LaTeX by draft.py; it was collated
%      word by word with an independent tesseract reading (dan, tessdata_best) of the same images (collate/check.py,
%      witness 'ocr'): 650 disagreements.  check.py's own rules settled 44.
%  (b) A THIRD reading voted: Project Runeberg's OCR of the same Google images (never proofread there), by context
%      (rules.py umn-*, 372 decisions), guarded so that a vote never makes a word the corpus does not attest; other
%      rules (the layer's glyph faults against an attested word, accents tesseract adds, the ø rule, the house thin space)
%      settled 70.  The vote was AUDITED at the image: 40 rule decisions sampled blind and 30 more suspects read at zooms:
%      7 Runeberg-for-the-draft votes were wrong (ABBYY and Runeberg sharing a fault: bg/og, kim/kun, bar/har,
%      Anstrengelser/Anstrængelser, Vist/vist, én/en, h Vad/hvad), all mended; the ~30 letter-changing votes for the draft
%      not looked at are tesseract noise against a real word.
%  (c) The rest were read at the image BLIND (round b1: 303 crops on 31 sheets, six readers, 62/62 planted controls right):
%      203 open checks, 23 places where Runeberg disputed both OCRs, 54 places where the henrikpontoppidan.dk e-text
%      differs (a FINDING AID only: someone else's transcription, never a vote -- it normalises the print's spelling, e.g.
%      Nietzsche for the 1889 print's Nietsche), and the audit; ~150 decided in the main session from the readers' lines,
%      ~25 at 300-dpi zooms.  OPEN 0.  An unattested-word sweep of the collated text (353 words in no other transcription)
%      was read through; its doubtful words were looked at in zooms (one misprint found: forkellige).
%  (d) STRUCTURE read at the image (reader K): the title pages, the article heads, the section numerals (two the draft
%      had lost), the verse, the divider; every FOOTNOTE read whole (reader N: 11 notes; compared word by word with the
%      collated notes: 0 word differences, the reader right on the closing quotes).
%  (e) LETTERSPACING (the print's emphasis, set \emph{} here) is MEASURED at the image, not read: two blind line rounds
%      (l1, 256 lines; l2, the 134 doubtful ones again) agreed on only 40 of 134 lines, so spacemeasure.py decides by the
%      white between the letters of each word in tesseract's boxes (bimodal: 1-4 px roman, 6-8 px spaced; threshold 5 px;
%      isolated words must have every gap wide).  Checked at zooms on 16 words where the measure and the e-text's <em>
%      differ: the measure right on 11; three of its false positives (heavy isolated words: Ævne, frem, Side) led to the
%      isolation rule; three misses mended by hand after the page check (SPACEFIX in spaceapply.py).
%  (f) SAMPLED PAGE CHECK (pagecheck.py, seed 20261015; two blind readers): 6 whole pages, one or two from each piece
%      (Religionsfrihed p. 451; Demokratisk Radikalisme pp. 852, 868; Gensvar pp. 132, 140; Epilog p. 343), 2202 words:
%      1 WORD error (Talemaader for the print's Talemaaler: a correction wrongly taken from the e-text finding aid),
%      1 PUNCT (a speck read as '), both mended; LETTERSPACING 15 differences: at 300-dpi zooms 5 transcription errors
%      (tre, Kopier twice, Genmæle, fortsatte; mended), 4 where the transcription is right, 6 not decidable at 300 dpi
%      (the ends of runs in wide-set lines).  Word error rate after mending 0 in 2202 (95% upper bound ~1.4 per 1000);
%      letterspacing remains the least certain feature: ~2-3 per 1000 words.
%  PRINTED AS IS (misprints kept, marked at the line):
%    p. 343: „Talemaaler“ for Talemaader?; the print Tale-/maaler (page check, zoom z14); the e-text has Talemaader
%  CONVENTIONS.  Page turns: % --- p. N --- and \opage{N} (the printed page); a word divided over a page turn is split with
%  %.  Footnotes as \footnote at their marks, numbered per page as printed (1)).  \emph{} = letterspacing; signatures in
%  small capitals as printed.  The running heads and the sheet signatures are not transcribed.  The ø rule (CLAUDE.md).
%  Quotation marks as printed: »...« and „...“ both occur, and the print sometimes doubles them (»»Zunächst).
% =====================================================================
\documentclass[12pt,a4paper,openany]{book}
\usepackage[utf8]{inputenc}
\usepackage[T1]{fontenc}
\usepackage{libertinus}
\usepackage{libertinust1math}
\usepackage[danish]{babel}
\usepackage[protrusion=true,expansion=true]{microtype}
\usepackage{setspace}
\usepackage[margin=1.3in]{geometry}
\usepackage{fancyhdr}
\usepackage{hyperref}
\hypersetup{pdftitle={Epilog}, pdfauthor={Harald Høffding}, hidelinks}
\pagestyle{fancy}
\fancyhf{}
\fancyhead[LE,RO]{\thepage}
\fancyhead[C]{\textit{H. Høffding: Epilog (1890)}}
\fancypagestyle{plain}{\fancyhf{}\fancyhead[LE,RO]{\thepage}}
\setstretch{1.3}
\usepackage{marginnote}
\usepackage{graphicx}
\DeclareUnicodeCharacter{0254}{\reflectbox{c}}   % ɔ (id est)
\renewcommand{\thefootnote}{\arabic{footnote})}   % the print numbers each page's notes 1), 2)
\newcommand{\opage}[1]{\marginnote{\footnotesize\textit{[#1]}}}
\newcommand{\sectionnum}[1]{\par\medskip\begin{center}#1\end{center}\smallskip}
\newcommand{\asterism}{\par\begin{center}*\quad*\\[-0.4em]*\end{center}}
\newcommand{\signature}[1]{\par\smallskip\hfill #1\hspace*{2em}\par}
\pdfstringdefDisableCommands{\renewcommand{\footnote}[2][]{}\renewcommand{\opage}[1]{}}

\begin{document}

\begin{titlepage}
\begin{center}
\vspace*{2em}
{\LARGE\bfseries \emph{TILSKUEREN.}}\\[1em]
{\large \emph{MAANEDSSKRIFT}}\\[0.5em]
{\small FOR}\\[0.5em]
LITERATUR, SAMFUNDSSPØRGSMAAL OG ALMENFATTELIGE\\
VIDENSKABELIGE SKILDRINGER.\\[1.5em]
{\small UDGIVET AF}\\[0.5em]
{\bfseries N. NEERGAARD.}\\[1.5em]
\emph{SYVENDE AARGANG.}\\
\vfill
KJØBENHAVN 1890.\\[0.4em]
{\bfseries P. G. PHILIPSENS FORLAG.}\\[0.3em]
{\small THIELES BOGTRYKKERI.}
\end{center}
\end{titlepage}
\clearpage

% --- p. 343 ---
\setcounter{footnote}{0}%
\opage{343}
\begin{center}{\large\bfseries Epilog.}\\[0.4em]\rule{3em}{0.4pt}\end{center}\medskip

Skøndt Dr. G. Brandes ved sin Vuggehistorie vel snarest 
har villet more Parterret, paa hvilket i det hele saa stor en Del
af hans Optræden i vor Strid har været anlagt, saa har den ogsaa
moret mig. Jeg kender tilstrækkeligt min egen Begrænsning til,
at jeg kan have mig selv til Bedste. Men Historien trænger til
et Supplement. Langt flere Gaver end de, den nævner, ere blevne
mig negtede. Jeg mangler en saadan Behændighedens Kunst, som
min Modstander i denne Strid har lagt for Dagen. Da jeg havde 
rettet min Kritik mod hans Artikel »Aristokratisk Radikalisme«, 
brækkede han resolut Spidsen af de store Ord og kraftige Talemaaler, % PRINTED AS IS (Talemaaler: for Talemaader?; the print Tale-/maaler (page check, zoom z14); the e-text has Talemaader)
hvoraf denne Artikel vrimlede, og søgte i en senere Artikel
(»Det store Menneske, Kulturens Kilde«) at fremstille sin Opfattelse
som den simpleste, uskyldigste Ting af Verden, medens han samtidigt
med utilbørlige og uretfærdige Angreb og Insinuationer gik
løs paa mig personligt. I sin tredje Artikel (»Duplik«) har han
til Dels villet bortforklare de Personligheder, hvoraf den anden 
Artikel vrimlede; men denne Bortforklaring er saa sindrigt indrettet, at den selv igen kan bortforklares. --- En saadan Kunst er 
mig negtet.

Skøndt Læserne sandsynligvis ere kede af den hele Strid,
der i Virkeligheden var udfægtet med de to første Artikler fra
hver Side, nødes jeg dog til i stor Korthed endnu en Gang at
belyse enkelte af Stridspunkterne, især for at faa Lejlighed til at
slaa fast, hvad det egentligt er, Striden har drejet sig om. Ti
desangaaende synes der selv hos velmenende og tænksomme
Læsere at være nogen Uklarhed. Jeg ønsker, at det skal staa
klart, hvad det i Virkeligheden er, jeg har villet til Livs.

\emph{Den rent personlige Side af Striden} skal jeg ikke
komme nærmere ind paa. Kun en uhistorisk Paastand maa jeg
% --- p. 344 ---
\setcounter{footnote}{0}%
\opage{344}berigtige. Dr. Brandes skriver, at de to Tidspunkter af hans Liv,
i hvilke jeg har været i Strid med ham, have været de, i hvilke
han havde ytret noget særligt upopulært. Her tager han storligt
Fejl. Under hans Strid mod den Kierkegaard-Nielsenske Lære om
Tro og Viden i Tredserne var han ikke upopulær, men blev tvært
imod støttet fra indflydelsesrige  Sider; han havde i det hele
Pressen for sig, og selv Teologerne saa med Sympati paa hans
Optræden. Man maa nemlig vel erindre, at den omtalte Lære om
Tro og Viden først og fremmest var rettet polemisk mod Teologien,
var et Forsøg paa at omstøde Teologien og dog hævde den
positivt religiøse Tro. I Halvfjærdserne førte Georg Brandes sin
egentlige, betydningsfulde Kamp, og i den havde han i det hele
min Sympati; jeg skrev i hans Tidsskrift og blev vist endog anset
for at staa ham nærmere end jeg i Virkeligheden gjorde. Den
Brandes, jeg nu er optraadt imod, er ikke denne Halvfjærdsernes,
men er Brandes fra Slutningen af Firserne.

Der synes at være en vis Rytme over mit Forhold til
Brandes efter de forskellige Decennier: i Tredsernes Modsætning,
i Halvfjærdsene Sympati, i Firserne atter Modsætning. Kan man
stole paa Rytmens Lov, er der jo gode Udsigter for Halvfemserne.
Og disse Udsigter vilde være endnu lysere, hvis han ikke havde
givet vor Strid en unødvendig personlig Karakter.

Kun ganske kort skal jeg komme ind paa \emph{de Steder af
Nietzsche, om hvis Fortolkning min Modstander og jeg
ere uenige}. --- Jeg har paany gennemlæst Stedet i »Genealogie
der Moral« p. 166---169 og maa --- om jeg saa nok saa meget bliver
beskyldt for Rethaveri --- fastholde min Opfattelse. Enkelte Udtalelser
af dette Stykke havde jeg i »Demokratisk Radikalisme«
anført for at vise, at Nietzsche var tilbøjelig til at mene, at den
strenge Prøvelse af Grundene for og imod en Sag, og den Tro
paa Sandheden, som ligger bag ved denne Strenghed, var en Rest
af Askese og Teologi. Den utrolige Dybsindighed, Dr. Brandes
finder i den tyske Forfatters Ord, er det mig umuligt at opdage.
Maaske mangler jeg Organet. De Læsere, som maatte interessere
sig for dette Spørgsmaal, maa selv efterlæse Stedet og overveje
det; men de maa paa Forhaand vide, at hvis de ikke komme til
samme Resultat som Dr. Brandes, ville de blive erklærede for
% --- p. 345 ---
\setcounter{footnote}{0}%
\opage{345}mere eller mindre bornerede. Om man anerkender Nietzsche eller
ikke, synes at skulle gøres til en Prøvesten for Aanderne ved
denne fin du siécle. ---

I sin Iver for at forsvare Nietzsche mod enhver Kritik er
Dr. Brandes kommen til at frakende ham Sammenhæng i Opfattelsen
paa et Punkt, hvor den virkeligt findes hos ham. I mit
Gensvar havde jeg sagt, at N. hævder, at der gives Følelser, som
ere bestemte blot ved Aktivitet, ikke ved noget Forhold til ydre
Betingelser, og som Bevis herfor havde jeg henvist til hans
Polemik mod den moderne Biologi, hvor han hævder »Livsviljens«
eller »Magtviljens« primære Karakter, »de spontane,
angribende, overgribende Kræfters« Uafhængighed af den Tillæmpelse,
paa hvilken Udviklingslæren lægger saa stor Vægt. (Geneal.
der Moral. p. 70). Hertil bemærker Dr. Brandes: »Ikke med en
Stavelse taler Nietzsche her om Følelser, ikke et Øjeblik sigter
han til den aristokratiske Selvfølelse«. Dr. Brandes, som særligt
paa dette Punkt finder Beviset for, at jeg mangler »det Gran af
Salt, hvormed det gælder at forstaa«, har her vistnok selv haft
et eller andet Gran i Øjet, der har hindret ham i at se, hvad jeg
ogsaa havde bemærket, at den hele lille Exkurs om den moderne
Biologi ikke er uden Sammenhæng med den nærmest foregaaende
Omtale af aktive og reaktive Følelser. Nietzsche vil forsvare Begrebet
Aktivitet (saaledes som han forstaar det) i Biologien, ligesom
han har hævdet det i Følelsens Psykologi, og Udtrykkene
pege derfor uafladeligt tilbage til det foregaaende Afsnit. (Smlgn.
f. Eks. »die spontanen, angreifenden, übergreifenden Kräfte« p. 70
med »der aktive, angreifende, übergreifende Mensch« p. 64). Deri
har N. ogsaa Ret, at naar der maatte anerkendes en »Aktivitet«
i biologisk Henseende, saa maatte der sikkert dertil svare en vis
Følelse, ligesom til Reaktiviteten. Naar Dr. Brandes oversætter
»Lebenswille« ved »Livsattraa«, bliver dette saa meget des
klarere. --- Med fuld Ret har jeg da kunnet benytte N.s Biologi
til at fortolke hans Psykologi.

Der er en tydelig Bestræbelse hos Dr. Brandes for at tillægge
mig saa bornerede Opfattelser som muligt. Saaledes mener
han, at naar jeg har sagt, at Nietzsche er Poet eller Retor, ikke
Filosof, gaar jeg ud fra udvortes, døde Rubriker, og at jeg har
% --- p. 346 ---
\setcounter{footnote}{0}%
\opage{346}en saa filistrøs \emph{Opfattelse af Filosofien}, at jeg kun kender
Synspunkterne: rigtigt-galt, og ikke har Blik for den Betydning, en
Tankegang kan have, naar den kommer fra en ejendommelig
Personlighed, selv om den ikke er »rigtig«. Intet strider mere
mod min hele Opfattelse af Filosofien og dens Historie end den
 Opfattelse, som her tillægges mig. Det er min Overbevisning, at
 i Filosofien, som hvad saa mange af dens Problemer angaar bevæger sig paa den menneskelige Erkendelses Grænseomraader,
maa nødvendigvis Forskerens Personlighed komme til at træde
mere frem end i andre Videnskaber. Erfaring viser da ogsaa, at
naar man indlader sig paa en Debat om filosofiske Æmner, maa
man give sin Personlighed til Pris, især naar man har en Modstander,
i hvis Kram det særligt passer at drage Striden ind paa
det personlige Omraade, maaske for derved at aflede Opmærksomheden
fra andre Sider. Jeg opfatter Filosofiens Historie som en stor,
fremskridende Dialog, i hvilken de forskellige Tankebevægelser,
som Livets og Erfaringens Problemer vække,' udfolde sig under
indbyrdes Kamp. Interessen ved de forskellige filosofiske Fænomener
beror paa den Energi og Konsekvens, med hvilken
Tankerne opstilles og gennemføres, og Interessen ved Filosofiens
Historie paa den Udviklingsproces, Tankerne maa gennemgaa
for at blive lutrede og prøvede. Der er Tanker og Problemløsninger,
hvis Tid synes at være omme, og som næppe en Gang
ville kunne hævde deres Plads i Mulighedernes Verden. Der er
andre Tanker, som synes at kunne udformes til at blive faststaaende
Sandheder. Der finder stadige Strømninger Sted mellem
Mulighedernes, Virkelighedens og Umulighedens Verden, og der er
en egen Interesse ved at følge deres Løb.

Men ligesom den enkelte Tanke maa undergaa en Lutring
og en Prøvelse, saaledes maa ogsaa den tænkende Personlighed
det. Ogsaa her er der en Vekselvirkning. Den filosofiske Tænkning
er ikke blot Udslag for en personlig Trang; den stiller ogsaa
sine bestemte Fordringer til den filosoferende Personlighed. Den
kræver Underkastelse og Resignation under Tænkningens Love,
Mod til »at forbyde sig sit Nej lige saa strængt som sit Ja«,
Ævne til at tilbagetrænge Stemninger, som ville hindre Tankens
faste Gang. En Forfatter kan være af fremragende Betydning ved
sin Ævne til at udtale stærke Stemninger, uden at besidde denne
Tankefasthed. Her tror jeg der savnes overordentligt meget hos
\emph{Nietzsche}, og jeg har derfor bebrejdet Dr. Brandes, at han har
% --- p. 347 ---
\setcounter{footnote}{0}%
\opage{347}indført ham som en stor Filosof og gentaget hans haanende Omtale
af redelige og energiske Forskere, hvis Skobaand Nietzsche
(naar han skal betragtes som Filosof) ikke er værdig at løse. Min
Optræden var ikke mindst dikteret af Pietet overfor Tænkere,
som vi alle skylde meget, og hvis Betydning ikke gendrives ved
romantiske Rodomontader. Jeg har ikke villet gøre mig selv stor,
men jeg har villet værge om det virkeligt store og sande.

Det er Mangelen af klar og virkelig Begrundelse af Paastandene
hos Nietzsche, som gør det saa vanskeligt at diskutere
om ham. Paa de mest afgørende Punkter findes der Modsigelser
hos ham. Hvis man nærmest opfatter ham som Digter, gør dette
intet, ti Digteren personificerer saa at sige sine Stemninger og
stiller dem mod hverandre som Personer i et Drama. Men af
en Tænker fordrer man dog i det mindste en Ledetraad til at
forlige Modsætningerne. --- Dr. Brandes har paavist Steder, hvor
N. omtaler Forskellen mellem Moralens historiske Opstaaen og
dens Gyldighed. Men alligevel behandler N. faktisk sin Teori om
den efter hans Mening herskende Morals Opstaaen som en tilstrækkelig
Kritik af denne Moral, og han meddeler ikke, hvilken
Maalestok han selv anvender ved sin Forkastelse af alle Værdier
(Entwerthung aller Werthe). Man kan ikke forkaste Værdi uden
i Kraft af andre Værdier; men hvorledes disses Gyldighed begrundes,
faa vi intet at vide om; vi afspises med historisk-profetiske
Udtalelser. Velfærdsprincipet forkaster N.; men hvorledes
han vil bære sig ad med at ræsonnere om Værdier (ɔ: Goder)
uden at antage det i en eller anden Form, med en eller anden
Modifikation, er ikke let at se, skøndt det er konsekvent af ham
at forkaste det, naar han gaar ud fra den Anskuelse, han har
udviklet i »Genealogie der Moral« --- men som rigtigt nok, som
vi skulle se, aldeles strider mod den, han i et andet Skrift lejlighedsvis
udtaler.

Overfor den engelske Biologi og dens strengt videnskabelige
Lære om Udviklingen af Livsformerne under Vekselvirkningen
med Omverdenen tager N. sin Tilflugt til det gamle mystiske Begreb
om en formende Drift, et Begreb, som aldeles intet forklarer,
ikke mere end om man vilde udlede Livet af en særegen »Livskraft«.
Det er atter Romantik. Og det er tillige Uretfærdighed
mod den moderne Biologi, som ikke negter, at al Udvikling paa
ethvert Punkt forudsætter noget, som kan udvikles, en given
Natur, der kan udfoldes videre gennem Vekselvirkning med Om%
% --- p. 348 ---
\setcounter{footnote}{0}%
\opage{348}verdenen. Særligt har Darwin stedse med klare Ord hævdet, at
den naturlige Udvælgelse kun kan øve sit Værk, naar der paa en
eller anden Maade er fremkommet Variationer. Hvorledes disse
Variationer opstaa, deri saa Darwin en stor Gaade og en foreløbig
Grænse for Forskningen; men han vilde sikkert ikke have
følt sig hjulpen ved Appellen til en formende Kraft, der kun er
en mystisk Tyskers Forsøg paa at restaurere den romantiske
Naturfilosofi. Vi have dog nu lært at skelne mellem hvad vi forstaa,
og hvad vi ikke forstaa, og ikke tage klingende Ord for Forklaringer
af det uforstaaede.

Endnu et Eksempel til Karakteristik af Nietzsche som
Tænker. I hans filosofiske Hovedskrift er det en Hovedtanke, som
atter og atter varieres, at der er stor Fare for, at de syge skulle
smitte de sunde og berøve dem deres Kraft og deres Lykke ved
at berøve dem Livsmodet. Han hævder Sygeligheden som det
normale og fordrer de faa sunde og lykkelige isolerede, for at
de ikke ogsaa skulle gaa til Grunde. (Geneal. der Moral, p. 132 ff.).
Derimod i et andet Skrift, som vel er nogle Aar ældre, men optryktes
paany i samme Aar som hint udkom, finder han det
netop karakteristisk for vore Dage, at de fleste af os ikke kende
Nød og Lidelse, hverken aandelig eller legemlig. Heraf (»aus der
allgemeinen Ungeübtheit in Schmerz beiderlei Gestalt und einer
gewissen Seltenheit des Anblicks eines Leidenden«) udleder han
et forøget Had til Smerten og Mangelen paa Ævne til endog at
udholde Tanken om Smerte. (Die fröhliche Wissenschaft. Neue
Ausgabe. p. 74 ff.). Hvad der her erklæres for et sjældent Syn,
er hist saa hyppigt, at det gør Europa til et Hospital og medfører
en social Fare. Det er aabenbart modsatte Stemninger, der
her besjæle ham, og hver Stemning udarbejdes til en Teori. Kun
én af disse \emph{Teorier} kan dog være sand, selv om hver af de to
\emph{Stemninger} have deres Berettigelse. Begge Stemninger ere
interessante, og deres Modsætning udelukker langtfra, at N. er en
fremragende Forfatter. Men til Fordel for hans Filosofi taler
denne Modstrid ikke. ---

Dog --- Uenighed om Nietzsches Betydning har ikke været
Stridens egentlige Aarsag.

Heller ikke \emph{de store Menneskers Betydning} har
egentligt været Stridens Aarsag. Jeg kan efter min Synsmaade
% --- p. 349 ---
\setcounter{footnote}{0}%
\opage{349}give dem alt hvad der kan fordres for dem, og jeg lægger den
aller største Vægt paa deres Uundværlighed. Ligesom det er en
Grundtanke i min Etik (se særligt Afsnittene om Kulturen og
om Staten), at alt nyt og værdifuldt har sit Udspring i et enkelt
Menneskes Sind, ud fra hvilket det træder frem for at kæmpe
sin Kamp for at bestaa og herske i Verden, saaledes havde jeg
allerede i et ældre Skrift (Om Grundlaget for den humane Etik.
1876. p. 81 f.) udtalt: »Historiens store Skikkelser, i hvilke de
menneskelige Opgaver og Formaal fremtræde ligesom personificerede,
staa som Forbilleder, som Typer for, hvad det menneskelige
Liv kan være«, og jeg havde anført Longfellows smukke
Linjer (af hans Psalm of Life):

\begin{center}\small Lives of great men all remind us,\\
We can make our lives sublime.\end{center}

Der vil heller ikke kunne paavises noget i mine seneste
Udtalelser, af hvilket der skulde kunne udledes hvad Dr. Brandes
skriver om mig i sin Duplik: »Det ønskelige var for ham, at\dots{}
Livet gled stilfærdigt af Sted som en douce Strøm med en lille
fredelig Gennemsnitslykke for Hvermand«. Dette følger ingenlunde
af min Ytring om, at det ofte var ønskeligt, om de store Mennesker
havde kunnet undværes. Jeg tænkte derved særligt paa de ulykkelige
Tider, Opløsningens Tider (som i Følge Nietzsche især ere dem, i
hvilke de store og fornemme blive til). Naar man kalder de
store Mennesker Historiens Maal, saa maa dog vel deri ligge, at
det altid er ønskeligt at skaffe dem Betingelserne for at optræde.
Men hvem vilde styrte et Folk i Ulykke, for at en stor Mand
kunde faa Lejlighed til at virke? Et helt andet er det, hvad man
bag efter kan sige. Naturligvis har Jeanne d'Arc's Skikkelse været
og er stadigt et stort Gode for Frankrig. Jeg har selv gjort opmærksom
paa, hvorledes Modsætningerne: Fald og Storhed i
Virkeligheden ere sammenkedede. --- Jeg kunde dog godt ogsaa
udstrække Betragtningen til andre Tilfælde. Dersom vi kunde
gøre Livet i \emph{alle} saa kraftigt, rigt og ophøjet, at de enkelte
særligt saa kaldte store ikke vare nødvendige som Hjælpere eller
Forbilleder, skulde det saa ikke være det bedste? Og er dette i
en vis Forstand ikke det egentlige Maal, som leder os i vor
Vurdering og Stræben i de enkelte Tilfælde?

Alle Mennesker ere baade Maal og Midler, og ingen af
disse to Synspunkter kan nogensinde falde helt bort, eller: Idealet,
% --- p. 350 ---
\setcounter{footnote}{0}%
\opage{350}hvorefter vi etisk set stræbe, er det, at de to Synspunkter ingensinde
blive skilte ad. Den Fortolkning, jeg har givet af Velfærdsprincipet,
gaar netop ud herpaa. (Se min Etik. p. 31. 217 f.
o. fl. St.). --- Om de store Mennesker har jeg saa i mine Artikler
sagt, at de ere store ved at være de store Midler, d. v. s. ved at
der fra dem udgaar betydningsfulde Virkninger, som højne det
menneskelige Liv og føre det fremad. Kun fordi de-ere de store
Midler, blive de ogsaa store Formaal, idet det paa Grund af det
store Bidrag, hvormed de forøge Slægtens Kapital, bliver af afgørende
Betydning for det menneskelige Liv, at saadanne Mennesker
opstaa. Men deres Opstaaen skjuler sig des værre i det dunkle, og
Vægten maa derfor lægges paa at tilvejebringe Betingelser for,
at de kunne blive anerkendte og komme til at virke. Dette sker
nu netop ved, at den hidtil vundne Kultur udbredes til saa mange
som muligt. Jo større Udsæd, des større Mulighed for særligt
heldige Frembringelser, og des større Mulighed for Forstaaelse af
det fremragende. Maalestokken for den elementære Kulturs Dygtighed
sætter jeg netop deri, om den frembringer en voksende
Ævne til at anerkende det store, hvor det virkeligt findes. Forholdet
mellem den højere og den elementære Kultur bør være
som mellem Digteren og Kritikeren, mellem de produktive og de
receptive Talenter. Ævnen til Forstaaelse kan stedse vokse hos
de mange, selv om Ævnen til at frembringe stedse vil være indskrænket
til faa (som dog stedse opstaa rundt om i Kredsen af
de mange). Der kan være et aandeligt Fællesskab til Stede, selv
om Ævnerne ere forskellige, og det er den vigtigste Opgave at
udvikle dette Fællesskab videre. Ti et saadant Fællesskab kan
med Rette siges at være Maalet.

Velfærdsmoralen er for mig ikke en dogmatisk Formel,
men en Hypotese, med hvilken jeg arbejder, fordi jeg ingen bedre
kender, og fordi enhver anden, jeg kender, indirekte forudsætter
den. Den stiller i Virkeligheden ethvert etisk Problem, saaledes
ogsaa Problemet om de store Mænds Betydning. Jeg er ikke
blind for, at de demokratiske Ideer i Teori og Praksis kunne
misforstaas, og ofte misforstaas, som om det nu var ligegyldigt,
om der opstod store, fremragende Aander eller ikke. Jeg har
selv haft Lejlighed til, at udtale i et Foredrag, at Demokratiet
vilde misforstaa sig selv, hvis det vilde mene, at det ikke behøvede
Mænd, der vare et Hoved højere end andre Mænd. Og
for længe siden har Stuart Mill set Problemet og behandlet det i
% --- p. 351 ---
\setcounter{footnote}{0}%
\opage{351}sin udødelige Bog »Om Friheden« med den store, sandhedskærlige
Tænkers hele, virkelige Overlegenhed. Han var ikke af dem, der
ikke kunde betone en for Tiden gavnlig Sandhed, uden at bruge
saa store Ord, at alle andre Sandheder stødtes til Side. Faren er
der, og hvis det blot havde drejet sig om »zeitgemässe Betrachtungen«,
om at indskærpe noget, der nu til Dags ikke sjældent
blev overset, vilde det have haft min fulde Sympati. Men ingen
skal faa mig til at tro, at Udtalelser af den Art, som Artiklen
»Aristokratisk Radikalisme« vrimlede af, vare nødvendige dertil,
eller egnede sig dertil.

Den Fare, som er til Stede, kan, efter min Overbevisning,
kun besværges ved ærligt og udholdende Arbejde indenfor
\emph{Demokratiet selv}. Jeg har den faste Tro, som ikke er ubegrundet
ved Erfaring (ti jeg er ikke den Stuelærde, Dr. Brandes
har villet gøre mig til, skøndt mine Salonerfaringer ganske vist ere
tarvelige), --- at der er meget at udrette, dersom vi kunne faa
Mennesker med Ævne og Vilje til at arbejde for »Massen«. Det
vil vise sig, at denne i Virkeligheden bestaar af enkelte. Konsekvensen af Principet om »den enkelte« er ikke den at skille sig
fra Massen paa fornem Vis; den eneste Maade, paa hvilken man
kan faa Bugt med Massen, er at opløse den i enkelte; vil man
øve Tryk imod den ovenfra, vil man blot gøre den saa meget
mere kompakt. Et Aristokrati vil aldrig faa Bugt med Massen,
men netop styrke den. --- Alle virkeligt etiske Kulturopgaver og
Reformideer hvile paa en Tankegang som denne.

Nietzsche er endnu langt fjærnere fra denne Tankegang end
man faar Indtrykket af gennem Fremstillingen i »Aristokratisk
Radikalisme«. Han fordømmer baade Protestantismen og den
franske Revolution, fordi de vare Opstande mod en fornem, fint
dannet og mistænksom Menneskehed og førte til Sejer for det
samme Princip, som forhen havde givet Kristendommen Sejer over
den klassiske Oldtid. Reformationen var efter hans Opfattelse en
aandelig Bondeopstand, en Pøbelbevægelse, og Luther en Bondelømmel,
som ærgrede sig over Romerkirkens Fornemhed og
aabnede det Allerhelligste for alle! (Die fröhliche Wissenschaft.
p. 274. 292. Geneal. der Moral. p. 36. 162). --- Den virkelige
aristokratiske Radikalisme maa drage saadanne Konsekvenser.
Den demokratiske Radikalisme frygter derimod ikke for at aabne
det Allerhelligste for de mange. Den véd, at man ikke behøver
at holde de ikke virkeligt indviede borte fra det hellige; de
% --- p. 352 ---
\setcounter{footnote}{0}%
\opage{352}\emph{kunne} slet ikke faa Adgang til det, selv om den »aabnes«. Men
den véd tillige, at det som først hvidskes i Ørerne, senere skal
forkyndes paa Tagene, for at Slægten kan leve paa det, indtil
atter nye Helligdomme stiftes i de enkeltes Sind, hos hvem de
nye livsbærende Tanker fødes. Her er et stadigt Kredsløb.

Naar Nietzsche beundrer Napoleon, da er det ikke paa
Carlyles og Brandes' æstetiserende Maade. Han priser Napoleon,
fordi man kan takke ham for, »at der nu formodentligt vil følge
et Par krigerske Aarhundreder efter hinanden, som ikke have deres
Lige i Historien, kort sagt, at vi ere indtraadte i Krigenes klassiske
Tidsalder«, og fordi han »i de moderne Ideer og i hele Civilisationen
saa ligesom en personlig Fjende og ved dette Fjendskab
viste sig som en af de største Fortsættelser af Renaissancen«.
(Die fröhliche Wissenschaft. p. 301 f.). Man ser her tydeligt den
Sammenhæng med den tyske Militarisme, som jeg allerede i mit
Gensvar (p. 139) pegede hen paa til Forstaaelse af Nietzsche.
Man maa lade denne, at han gør Alvor af det »fornemme« Ideal og
tramper talløse Skarer ned, for at en enkelt kan komme til at
vise sig som en »Mand«. Han er derfor konsekvent, naar han
afviser Velfærdsprincipet. Men hvorledes man kan hylde hans
Grundtanker og dog fastholde Velfærdsprincipet, det er mig endnu
uforstaaeligt. Og om dette Punkt yder ikke heller Dr. Brandes'
Duplik nogen Oplysning.

Naar det hverken var om Nietzsche eller om Betimeligheden
af at hævde de store Mænds Betydning, at Striden opstod, hvad
var da dens Æmne?

En Mand, som kun havde læst Dr. Brandes' anden Artikel
og mit »Gensvar«, ikke vore første Artikler, kunde ikke begribe,
hvad vi egentligt vare uenige om. Dette forstod jeg meget vel.
Ti det, som enhver maa gøre for at danne sig en rigtig Forestilling
om Stridens Aarsag og for at fælde en rigtig Dom om
min Optræden i denne Sag, maa være at sammenligne Dr. Brandes'
første Artikel (»Aristokr. Radikalisme«) med hans anden Artikel
(»Det store Menneske, Kulturens Kilde«), som han skrev efter mit
Angreb.

\emph{Dersom Dr. Brandes fra først af havde udtalt sig
om de store Mennesker, som han har gjort i sin anden}
% --- p. 353 ---
\setcounter{footnote}{0}%
\opage{353}\emph{Artikel og i sin Duplik, vilde jeg ikke have angrebet}
ham. Men som han i den første Artikel udtalte sig gennem
Aand og Tone, saa vel som gennem flere enkelte Udtryk, kunde
jeg ikke forstaa hans Opfattelse anderledes end som \emph{en absurd
Forgudelse af enkelte »store« paa Bekostning af det,
der er Historiens og særligt den nyere Tids Ære: Udvidelsen
af Menneskelighedsfølelsen til i Principet at
omfatte alle}. Det var denne Art Paastande, som jeg mente,
behøvede en bestemt og alvorlig Indsigelse og Kritik. Den, som
nu for Alvor vil dømme Dr. Brandes og mig imellem, maa undersøge,
om det ikke forholder sig, som jeg har sagt, nemlig, at
hans anden Artikel betegner et afgjort, 'et tydeligt Tilbagetog i
Forhold til den første.

I mit Gensvar (p. 140 f. og 148—152) havde jeg søgt at
vise, at der havde fundet et saadant Tilbagetog Sted, og at Dr.
Brandes ikke uden Selvmodsigelse kunde fastholde sin Tilslutning
til Nietzsche og dog være uenig med denne i hans Forkastelse af
Velfærdsmoralen. Paa disse Punkter maa jeg bede enhver fæste
sin Opmærksomhed, hvis han vil have en Mening om vor Strid.

Jeg havde ventet, at Dr. Brandes i sin Duplik var gaaet
nærmere ind paa disse Punkter og vilde have søgt at værge sig over
for de nævnte i mit Gensvar opstillede Paastande. Det har han ikke
gjort, skøndt det havde været rimeligere at anvende Pladsen dertil
end til at sige mig Ubehageligheder og give sin Vittighed Luft
om den mærkelige Ting, at jeg er Professor. (Herre Gud, vi
kunne jo ikke alle være Profeter). Naar han intet har haft at
sige om hine Punkter i sin Duplik, har jeg i min Epilog naturligvis
heller ikke Grund til at sige noget derom.

Dr. Brandes siger, at jeg har kæmpet mod Vejrmøller. Jeg
kan godt gaa ind herpaa, kun at jeg ikke tog disse Vejrmøller for
Riddere, men netop tog dem for hvad de vare. Artiklen om »Aristokratisk
Radikalisme« var fuld af klaprende Vejrmøller, og der blev
tillige i den gjort en saadan Blæst, at der opstod en Susen, som
nok kunde forvirre svage Hjærner og faa dem til at tro, at en ny
Æra var i Anmarsch. Hvad jeg vilde, var blot at faa disse Vejrmøller
jævnede med Jorden, for at de ikke længere skulde vække 
Ørhed og Forvirrelse. Og jeg kan ikke se andet, end at jeg har
opnaaet hvad jeg vilde.
\signature{\textsc{Harald Høffding}.}

\end{document}
